\documentclass[11pt,a4paper]{book}
\usepackage[margin=0.75in]{geometry}
\usepackage[T1]{fontenc}
\usepackage{lmodern}
\usepackage{enumitem}
\usepackage{xcolor}
\usepackage{hyperref}
\usepackage{longtable}
\usepackage{booktabs}
\usepackage{fancyhdr}
\hypersetup{colorlinks=true,linkcolor=blue,urlcolor=blue}
\setlist[enumerate]{itemsep=2pt,topsep=3pt}
\pagestyle{fancy}
\fancyhead[LE,RO]{JKSSB Computer Awareness MCQs}
\fancyhead[RE,LO]{Original chapterwise question bank}
\fancyfoot[C]{\thepage}
\newcommand{\source}[1]{\par\smallskip\textit{Source: #1}}
\newcommand{\answer}[2]{\par\textbf{Answer: #1.} #2}
\begin{document}
\frontmatter
\begin{titlepage}\centering\vspace*{2cm}{\Huge\bfseries JKSSB Computer Awareness\\[0.3em]MCQ Question Bank\par}\vspace{1cm}{\Large Complete syllabus coverage\par}\vspace{1cm}{\large 200 original MCQs with sources, answers and explanations\par}\vfill{\large Prepared in LaTeX\\October 2026\par}\end{titlepage}
\chapter*{How to use this book}
This is an original practice book built from the supplied ten-chapter syllabus. It follows the chapterwise breadth visible in Sanfoundry collections and the objective-test emphasis of official JKSSB syllabus and question-paper materials. Questions are newly written; they are not copied verbatim. Each chapter contains 20 questions, including at least one numerical/application item (5\% module minimum).
\textbf{Source labels.} ``Sanfoundry identifies the topic family used for coverage calibration; ``JKSSB official identifies the official syllabus/question-paper pattern; ``original synthesis means the question and explanation were authored for this book. Always check the latest JKSSB notification and official answer keys because version-sensitive items can change.
\chapter*{Source register}
\begin{itemize}
\item JKSSB official previous-question-paper archive: \url{https://jkssb.nic.in/QP.html}
\item JKSSB Junior Assistant/Stenographer syllabus notice (23 December 2025): \url{https://jkssb.nic.in/Pdf/Syllabus_JrAsstt_Steno_23122025.pdf}
\item Sanfoundry Computer Fundamentals chapterwise MCQs: \url{https://www.sanfoundry.com/1000-computer-fundamentals-questions-answers/}
\item Sanfoundry Operating System chapterwise MCQs: \url{https://www.sanfoundry.com/operating-system-questions-answers/}
\item Sanfoundry Computer Network chapterwise MCQs: \url{https://www.sanfoundry.com/computer-network-questions-answers/}
\end{itemize}
\tableofcontents
\mainmatter
\chapter{Computer Fundamentals and Digital Foundations}
\textit{20 questions; numerical/application items are marked in the answer explanation.}
\section*{Question 1}
In the IPO cycle, which stage transforms raw data into meaningful information?
\begin{enumerate}[label=\Alph*.]
\item Output
\item Input
\item Processing
\item Storage
\end{enumerate}
\answer{C}{Processing. Processing applies instructions to input data.}
\source{Sanfoundry, Computer Fundamentals chapterwise MCQs; original synthesis}
\section*{Question 2}
Which characteristic means that a computer can perform repetitive work without becoming tired?
\begin{enumerate}[label=\Alph*.]
\item Automation
\item Versatility
\item Accuracy
\item Diligence
\end{enumerate}
\answer{D}{Diligence. Diligence refers to consistent performance without fatigue.}
\source{Sanfoundry, Computer Fundamentals chapterwise MCQs; original synthesis}
\section*{Question 3}
A computer that represents data using continuously varying physical quantities is a(n):
\begin{enumerate}[label=\Alph*.]
\item Analog computer
\item Digital computer
\item Hybrid computer
\item Mainframe
\end{enumerate}
\answer{A}{Analog computer. Analog computers process continuous values.}
\source{Sanfoundry, Computer Fundamentals chapterwise MCQs; original synthesis}
\section*{Question 4}
Which class of computer is generally associated with the highest processing capability?
\begin{enumerate}[label=\Alph*.]
\item Microcomputer
\item Supercomputer
\item Minicomputer
\item Workstation
\end{enumerate}
\answer{B}{Supercomputer. Supercomputers are designed for very high-performance computation.}
\source{Sanfoundry, Computer Fundamentals chapterwise MCQs; original synthesis}
\section*{Question 5}
How many bits are in one byte?
\begin{enumerate}[label=\Alph*.]
\item 16
\item 4
\item 8
\item 32
\end{enumerate}
\answer{C}{8. One byte consists of eight bits. Numerical/application item.}
\source{Sanfoundry, Computer Fundamentals chapterwise MCQs; original synthesis}
\section*{Question 6}
Which unit is equal to four bits?
\begin{enumerate}[label=\Alph*.]
\item Kilobyte
\item Byte
\item Word
\item Nibble
\end{enumerate}
\answer{D}{Nibble. A nibble is four binary bits.}
\source{Sanfoundry, Computer Fundamentals chapterwise MCQs; original synthesis}
\section*{Question 7}
Which is the correct binary representation of decimal 10?
\begin{enumerate}[label=\Alph*.]
\item 1010
\item 1001
\item 1100
\item 1110
\end{enumerate}
\answer{A}{1010. 8+2 = 10, so 1010 is correct. Numerical/application item.}
\source{Sanfoundry, Computer Fundamentals chapterwise MCQs; original synthesis}
\section*{Question 8}
Which code is designed to represent characters from many writing systems?
\begin{enumerate}[label=\Alph*.]
\item ASCII only
\item Unicode
\item BCD
\item EBCDIC only
\end{enumerate}
\answer{B}{Unicode. Unicode provides a universal character repertoire.}
\source{Sanfoundry, Computer Fundamentals chapterwise MCQs; original synthesis}
\section*{Question 9}
Which number system uses sixteen symbols, 0–9 and A–F?
\begin{enumerate}[label=\Alph*.]
\item Octal
\item Binary
\item Hexadecimal
\item Decimal
\end{enumerate}
\answer{C}{Hexadecimal. Hexadecimal is base 16.}
\source{Sanfoundry, Computer Fundamentals chapterwise MCQs; original synthesis}
\section*{Question 10}
Which gate produces 1 only when both inputs are 1?
\begin{enumerate}[label=\Alph*.]
\item XOR
\item OR
\item NOT
\item AND
\end{enumerate}
\answer{D}{AND. The AND truth table has output 1 only for 1 AND 1.}
\source{Sanfoundry, Computer Fundamentals chapterwise MCQs; original synthesis}
\section*{Question 11}
The fifth computer generation is commonly associated with:
\begin{enumerate}[label=\Alph*.]
\item Artificial intelligence
\item Vacuum tubes
\item Magnetic drums
\item Punch cards
\end{enumerate}
\answer{A}{Artificial intelligence. Fifth-generation descriptions commonly emphasize AI-oriented systems.}
\source{Sanfoundry, Computer Fundamentals chapterwise MCQs; original synthesis}
\section*{Question 12}
Which is an example of an embedded computer?
\begin{enumerate}[label=\Alph*.]
\item Desktop PC used for browsing
\item Controller in a washing machine
\item General-purpose server
\item Supercomputer
\end{enumerate}
\answer{B}{Controller in a washing machine. Embedded computers are dedicated to a larger device.}
\source{Sanfoundry, Computer Fundamentals chapterwise MCQs; original synthesis}
\section*{Question 13}
Which device is primarily used to convert a paper document into a digital image?
\begin{enumerate}[label=\Alph*.]
\item Projector
\item Plotter
\item Scanner
\item Speaker
\end{enumerate}
\answer{C}{Scanner. A scanner captures printed material as digital data.}
\source{Sanfoundry, Computer Fundamentals chapterwise MCQs; original synthesis}
\section*{Question 14}
A word in computing generally refers to:
\begin{enumerate}[label=\Alph*.]
\item A printer command
\item Exactly one bit
\item Exactly four bits
\item A processor-dependent group of bits
\end{enumerate}
\answer{D}{A processor-dependent group of bits. Word size depends on the processor architecture.}
\source{Sanfoundry, Computer Fundamentals chapterwise MCQs; original synthesis}
\section*{Question 15}
Which item is information rather than raw data?
\begin{enumerate}[label=\Alph*.]
\item A monthly sales report
\item Unsorted transaction entries
\item Unlabelled sensor readings
\item Individual keystrokes
\end{enumerate}
\answer{A}{A monthly sales report. A report is processed and organized data.}
\source{Sanfoundry, Computer Fundamentals chapterwise MCQs; original synthesis}
\section*{Question 16}
Which is a limitation of computers?
\begin{enumerate}[label=\Alph*.]
\item They cannot store data
\item They lack independent human judgment
\item They cannot repeat tasks
\item They are always inaccurate
\end{enumerate}
\answer{B}{They lack independent human judgment. Computers execute instructions but do not possess human judgment by themselves.}
\source{Sanfoundry, Computer Fundamentals chapterwise MCQs; original synthesis}
\section*{Question 17}
Which technology replaced vacuum tubes in the second generation?
\begin{enumerate}[label=\Alph*.]
\item Microprocessors
\item Integrated circuits
\item Transistors
\item Optical fibre
\end{enumerate}
\answer{C}{Transistors. Second-generation computers used transistors.}
\source{Sanfoundry, Computer Fundamentals chapterwise MCQs; original synthesis}
\section*{Question 18}
The RGB model is primarily used for:
\begin{enumerate}[label=\Alph*.]
\item Scheduling processes
\item Representing sound
\item Compressing text
\item Representing colour in displays
\end{enumerate}
\answer{D}{Representing colour in displays. Red, green and blue components form display colours.}
\source{Sanfoundry, Computer Fundamentals chapterwise MCQs; original synthesis}
\section*{Question 19}
If a storage device holds 2 KB using the binary convention, how many bytes is that?
\begin{enumerate}[label=\Alph*.]
\item 2048
\item 2000
\item 1024
\item 4096
\end{enumerate}
\answer{A}{2048. 2 x 1024 = 2048 bytes. Numerical/application item.}
\source{Sanfoundry, Computer Fundamentals chapterwise MCQs; original synthesis}
\section*{Question 20}
Which sequence best describes basic computer operation?
\begin{enumerate}[label=\Alph*.]
\item Output–input–storage–processing
\item Input–processing–output–storage
\item Storage–output–input–processing
\item Processing–storage–input–output
\end{enumerate}
\answer{B}{Input–processing–output–storage. The conventional IPO flow is input, processing and output, with storage supporting it.}
\source{Sanfoundry, Computer Fundamentals chapterwise MCQs; original synthesis}

\chapter{Computer Organization, CPU, Memory and Hardware}
\textit{20 questions; numerical/application items are marked in the answer explanation.}
\section*{Question 1}
Which CPU unit performs arithmetic and logical operations?
\begin{enumerate}[label=\Alph*.]
\item Bus Interface
\item Control Unit
\item Register Unit
\item ALU
\end{enumerate}
\answer{D}{ALU. The ALU performs arithmetic and logical operations.}
\source{Sanfoundry, Computer Fundamentals chapterwise MCQs; original synthesis}
\section*{Question 2}
The instruction cycle normally begins with:
\begin{enumerate}[label=\Alph*.]
\item Fetch
\item Execute
\item Store result
\item Format disk
\end{enumerate}
\answer{A}{Fetch. The processor fetches the next instruction from memory.}
\source{Sanfoundry, Computer Fundamentals chapterwise MCQs; original synthesis}
\section*{Question 3}
Which register normally contains the address of the next instruction?
\begin{enumerate}[label=\Alph*.]
\item Instruction Register
\item Program Counter
\item Accumulator
\item Memory Buffer Register
\end{enumerate}
\answer{B}{Program Counter. The Program Counter points to the next instruction.}
\source{Sanfoundry, Computer Fundamentals chapterwise MCQs; original synthesis}
\section*{Question 4}
Which memory is closest to the CPU and usually faster than RAM?
\begin{enumerate}[label=\Alph*.]
\item Magnetic tape
\item Optical disk
\item Cache
\item Cloud storage
\end{enumerate}
\answer{C}{Cache. Cache stores frequently used data near the processor.}
\source{Sanfoundry, Computer Fundamentals chapterwise MCQs; original synthesis}
\section*{Question 5}
Which memory is volatile?
\begin{enumerate}[label=\Alph*.]
\item Optical disk
\item ROM
\item Flash memory
\item RAM
\end{enumerate}
\answer{D}{RAM. RAM loses its contents when power is removed.}
\source{Sanfoundry, Computer Fundamentals chapterwise MCQs; original synthesis}
\section*{Question 6}
Which RAM type is commonly used for CPU cache because it is faster?
\begin{enumerate}[label=\Alph*.]
\item SRAM
\item DRAM
\item PROM
\item EPROM
\end{enumerate}
\answer{A}{SRAM. SRAM is faster and commonly used in cache.}
\source{Sanfoundry, Computer Fundamentals chapterwise MCQs; original synthesis}
\section*{Question 7}
Which memory technology can be electrically erased and reprogrammed?
\begin{enumerate}[label=\Alph*.]
\item PROM
\item EEPROM
\item Mask ROM
\item Magnetic tape
\end{enumerate}
\answer{B}{EEPROM. EEPROM is electrically erasable and programmable.}
\source{Sanfoundry, Computer Fundamentals chapterwise MCQs; original synthesis}
\section*{Question 8}
Virtual memory uses secondary storage to extend the apparent capacity of:
\begin{enumerate}[label=\Alph*.]
\item The monitor
\item The keyboard
\item Main memory
\item The ALU
\end{enumerate}
\answer{C}{Main memory. Virtual memory provides an extension of RAM using disk/SSD space.}
\source{Sanfoundry, Operating System chapterwise MCQs; original synthesis}
\section*{Question 9}
Which device normally provides persistent magnetic storage?
\begin{enumerate}[label=\Alph*.]
\item Cache
\item RAM
\item CPU register
\item HDD
\end{enumerate}
\answer{D}{HDD. A hard disk drive stores data persistently on magnetic media.}
\source{Sanfoundry, Computer Fundamentals chapterwise MCQs; original synthesis}
\section*{Question 10}
Which is generally faster for random access: SSD or HDD?
\begin{enumerate}[label=\Alph*.]
\item SSD
\item HDD
\item Magnetic tape
\item CD-ROM
\end{enumerate}
\answer{A}{SSD. SSDs have no moving heads and generally offer lower access latency.}
\source{Sanfoundry, Computer Fundamentals chapterwise MCQs; original synthesis}
\section*{Question 11}
The motherboard primarily:
\begin{enumerate}[label=\Alph*.]
\item Displays images
\item Connects major computer components
\item Prints documents
\item Scans barcodes
\end{enumerate}
\answer{B}{Connects major computer components. The motherboard provides sockets, buses and interconnections.}
\source{Sanfoundry, Computer Fundamentals chapterwise MCQs; original synthesis}
\section*{Question 12}
A bus in a computer is used to:
\begin{enumerate}[label=\Alph*.]
\item Store passwords only
\item Cool the CPU
\item Transfer data, addresses or control signals
\item Create pixels
\end{enumerate}
\answer{C}{Transfer data, addresses or control signals. Buses carry signals among components.}
\source{Sanfoundry, Computer Fundamentals chapterwise MCQs; original synthesis}
\section*{Question 13}
A multi-core processor contains:
\begin{enumerate}[label=\Alph*.]
\item Multiple operating systems necessarily
\item Several keyboards
\item Only one transistor
\item More than one processing core
\end{enumerate}
\answer{D}{More than one processing core. Each core can execute instruction streams subject to system design.}
\source{Sanfoundry, Computer Fundamentals chapterwise MCQs; original synthesis}
\section*{Question 14}
A GPU is specialized primarily for:
\begin{enumerate}[label=\Alph*.]
\item Parallel graphics and data processing
\item Printing text
\item Managing email addresses
\item Replacing all storage
\end{enumerate}
\answer{A}{Parallel graphics and data processing. GPUs contain many parallel processing units suited to graphics and data workloads.}
\source{Sanfoundry, Computer Fundamentals chapterwise MCQs; original synthesis}
\section*{Question 15}
Which architecture stores instructions and data in the same main memory?
\begin{enumerate}[label=\Alph*.]
\item Harvard-only architecture
\item Von Neumann architecture
\item Peer-to-peer architecture
\item Ring architecture
\end{enumerate}
\answer{B}{Von Neumann architecture. Von Neumann architecture uses a shared memory concept.}
\source{JKSSB official syllabus plus current-topic extension; original synthesis}
\section*{Question 16}
In immediate addressing, the operand is:
\begin{enumerate}[label=\Alph*.]
\item Only on magnetic tape
\item Always in a network server
\item Included in the instruction
\item Stored in a printer buffer
\end{enumerate}
\answer{C}{Included in the instruction. Immediate mode places the operand value inside the instruction.}
\source{JKSSB official syllabus plus current-topic extension; original synthesis}
\section*{Question 17}
If a cache access takes 2 ns and a main-memory access takes 50 ns, how much faster is main memory access time than cache access time?
\begin{enumerate}[label=\Alph*.]
\item 100 times
\item 2 times
\item 48 times
\item 25 times
\end{enumerate}
\answer{D}{25 times. 50 / 2 = 25. Numerical/application item.}
\source{JKSSB official syllabus plus current-topic extension; original synthesis}
\section*{Question 18}
Pipelining improves instruction throughput by:
\begin{enumerate}[label=\Alph*.]
\item Overlapping stages of multiple instructions
\item Removing the control unit
\item Making all instructions identical
\item Disabling registers
\end{enumerate}
\answer{A}{Overlapping stages of multiple instructions. Different instructions can occupy different pipeline stages simultaneously.}
\source{JKSSB official syllabus plus current-topic extension; original synthesis}
\section*{Question 19}
A data hazard occurs when:
\begin{enumerate}[label=\Alph*.]
\item A monitor has low brightness
\item An instruction depends on an unfinished result
\item A disk is formatted
\item A user changes a password
\end{enumerate}
\answer{B}{An instruction depends on an unfinished result. Pipeline hazards include dependencies between nearby instructions.}
\source{JKSSB official syllabus plus current-topic extension; original synthesis}
\section*{Question 20}
Which is normally non-volatile?
\begin{enumerate}[label=\Alph*.]
\item SRAM cache
\item DRAM
\item ROM
\item CPU register
\end{enumerate}
\answer{C}{ROM. ROM retains contents without power.}
\source{Sanfoundry, Computer Fundamentals chapterwise MCQs; original synthesis}

\chapter{Input, Output, Peripherals and Ports}
\textit{20 questions; numerical/application items are marked in the answer explanation.}
\section*{Question 1}
OCR is used to recognize:
\begin{enumerate}[label=\Alph*.]
\item Printed or typed characters
\item Marked answer bubbles
\item Magnetic ink only
\item Sound frequencies
\end{enumerate}
\answer{A}{Printed or typed characters. Optical Character Recognition converts printed characters into editable text.}
\source{Sanfoundry, Computer Fundamentals chapterwise MCQs; original synthesis}
\section*{Question 2}
OMR is commonly used to read:
\begin{enumerate}[label=\Alph*.]
\item Printed books
\item Marked forms
\item Bank cheques by ink
\item Audio recordings
\end{enumerate}
\answer{B}{Marked forms. Optical Mark Recognition detects filled marks.}
\source{Sanfoundry, Computer Fundamentals chapterwise MCQs; original synthesis}
\section*{Question 3}
MICR is associated particularly with:
\begin{enumerate}[label=\Alph*.]
\item Fingerprints
\item QR codes
\item Magnetic-ink characters on cheques
\item Video frames
\end{enumerate}
\answer{C}{Magnetic-ink characters on cheques. MICR reads characters printed with magnetic ink.}
\source{Sanfoundry, Computer Fundamentals chapterwise MCQs; original synthesis}
\section*{Question 4}
Which device is both an input and output device?
\begin{enumerate}[label=\Alph*.]
\item Plotter
\item Keyboard
\item Microphone
\item Touchscreen
\end{enumerate}
\answer{D}{Touchscreen. A touchscreen displays output and accepts touch input.}
\source{Sanfoundry, Computer Fundamentals chapterwise MCQs; original synthesis}
\section*{Question 5}
Which is an impact printer?
\begin{enumerate}[label=\Alph*.]
\item Dot-matrix printer
\item Laser printer
\item Inkjet printer
\item Thermal printer
\end{enumerate}
\answer{A}{Dot-matrix printer. Dot-matrix printers strike an ink ribbon.}
\source{Sanfoundry, Computer Fundamentals chapterwise MCQs; original synthesis}
\section*{Question 6}
Which printer is generally a non-impact page printer?
\begin{enumerate}[label=\Alph*.]
\item Dot-matrix printer
\item Laser printer
\item Daisy-wheel printer
\item Line printer
\end{enumerate}
\answer{B}{Laser printer. Laser printers form pages without mechanical striking.}
\source{Sanfoundry, Computer Fundamentals chapterwise MCQs; original synthesis}
\section*{Question 7}
A plotter is especially suitable for:
\begin{enumerate}[label=\Alph*.]
\item Scanning fingerprints
\item Playing audio
\item Engineering drawings
\item Sending email
\end{enumerate}
\answer{C}{Engineering drawings. Plotters produce large, precise line drawings.}
\source{Sanfoundry, Computer Fundamentals chapterwise MCQs; original synthesis}
\section*{Question 8}
Screen resolution is commonly expressed as:
\begin{enumerate}[label=\Alph*.]
\item Volts per inch
\item Bytes per second
\item Pages per minute
\item Width by height in pixels
\end{enumerate}
\answer{D}{Width by height in pixels. Resolution describes the number of pixels horizontally and vertically.}
\source{Sanfoundry, Computer Fundamentals chapterwise MCQs; original synthesis}
\section*{Question 9}
Refresh rate is measured in:
\begin{enumerate}[label=\Alph*.]
\item Hertz
\item Bytes
\item Dots per inch only
\item Pages per minute
\end{enumerate}
\answer{A}{Hertz. Refresh rate is cycles/frames per second, measured in hertz.}
\source{Sanfoundry, Computer Fundamentals chapterwise MCQs; original synthesis}
\section*{Question 10}
Which connector is commonly used for digital video and audio to a display?
\begin{enumerate}[label=\Alph*.]
\item PS/2
\item HDMI
\item RJ-11
\item SATA
\end{enumerate}
\answer{B}{HDMI. HDMI carries digital audio and video.}
\source{Sanfoundry, Computer Fundamentals chapterwise MCQs; original synthesis}
\section*{Question 11}
RJ-45 is commonly associated with:
\begin{enumerate}[label=\Alph*.]
\item SATA storage
\item Telephone handset audio
\item Ethernet networking
\item Monitor power
\end{enumerate}
\answer{C}{Ethernet networking. RJ-45 connectors are widely used for Ethernet cables.}
\source{Sanfoundry, Computer Fundamentals chapterwise MCQs; original synthesis}
\section*{Question 12}
USB-C describes primarily a:
\begin{enumerate}[label=\Alph*.]
\item Printer language
\item File system
\item Network protocol
\item Physical connector/reversible port form
\end{enumerate}
\answer{D}{Physical connector/reversible port form. USB-C is a connector form; the supported USB standard may vary.}
\source{JKSSB official syllabus plus current-topic extension; original synthesis}
\section*{Question 13}
Which device converts digital data into a paper document?
\begin{enumerate}[label=\Alph*.]
\item Printer
\item Scanner
\item Webcam
\item Microphone
\end{enumerate}
\answer{A}{Printer. A printer produces hard copy output.}
\source{Sanfoundry, Computer Fundamentals chapterwise MCQs; original synthesis}
\section*{Question 14}
Which device is designed to capture live video?
\begin{enumerate}[label=\Alph*.]
\item Plotter
\item Webcam
\item Projector
\item Speaker
\end{enumerate}
\answer{B}{Webcam. A webcam captures video input.}
\source{Sanfoundry, Computer Fundamentals chapterwise MCQs; original synthesis}
\section*{Question 15}
A barcode reader is primarily an:
\begin{enumerate}[label=\Alph*.]
\item Storage device
\item Output device
\item Input device
\item Processing unit
\end{enumerate}
\answer{C}{Input device. It reads encoded information into the computer.}
\source{Sanfoundry, Computer Fundamentals chapterwise MCQs; original synthesis}
\section*{Question 16}
A pixel is best described as:
\begin{enumerate}[label=\Alph*.]
\item A CPU instruction
\item A printer cartridge
\item A network packet
\item The smallest addressable element of a digital image/display
\end{enumerate}
\answer{D}{The smallest addressable element of a digital image/display. Pixel means picture element.}
\source{Sanfoundry, Computer Fundamentals chapterwise MCQs; original synthesis}
\section*{Question 17}
If a display has 1920 horizontal pixels and 1080 vertical pixels, its total pixel count is:
\begin{enumerate}[label=\Alph*.]
\item 2,073,600
\item 2,880
\item 3,000,000
\item 1,920,000
\end{enumerate}
\answer{A}{2,073,600. 1920 x 1080 = 2,073,600 pixels. Numerical/application item.}
\source{Sanfoundry, Computer Fundamentals chapterwise MCQs; original synthesis}
\section*{Question 18}
Which port is most closely associated with internal SATA drives?
\begin{enumerate}[label=\Alph*.]
\item HDMI
\item SATA
\item VGA
\item RJ-45
\end{enumerate}
\answer{B}{SATA. SATA is a storage interface.}
\source{Sanfoundry, Computer Fundamentals chapterwise MCQs; original synthesis}
\section*{Question 19}
Soft copy refers to:
\begin{enumerate}[label=\Alph*.]
\item A punched card only
\item Printed paper
\item Electronic representation on a screen or file
\item A physical cable
\end{enumerate}
\answer{C}{Electronic representation on a screen or file. Soft copy is not a physical printed output.}
\source{Sanfoundry, Computer Fundamentals chapterwise MCQs; original synthesis}
\section*{Question 20}
Which process converts continuous audio into digital samples?
\begin{enumerate}[label=\Alph*.]
\item Indexing
\item Plotting
\item Defragmentation
\item Sampling
\end{enumerate}
\answer{D}{Sampling. Sampling measures a continuous signal at discrete intervals.}
\source{JKSSB official syllabus plus current-topic extension; original synthesis}

\chapter{Software, Operating Systems and Windows}
\textit{20 questions; numerical/application items are marked in the answer explanation.}
\section*{Question 1}
Which software manages hardware resources and provides services to applications?
\begin{enumerate}[label=\Alph*.]
\item Spreadsheet
\item Operating system
\item Compiler only
\item Web page
\end{enumerate}
\answer{B}{Operating system. The operating system manages processes, files, memory and devices.}
\source{Sanfoundry, Operating System chapterwise MCQs; original synthesis}
\section*{Question 2}
The core component of an operating system is the:
\begin{enumerate}[label=\Alph*.]
\item Spreadsheet
\item Toolbar
\item Kernel
\item Browser tab
\end{enumerate}
\answer{C}{Kernel. The kernel manages core system operations.}
\source{Sanfoundry, Operating System chapterwise MCQs; original synthesis}
\section*{Question 3}
GUI stands for:
\begin{enumerate}[label=\Alph*.]
\item Graphic Unit Instruction
\item General Utility Index
\item Global User Internet
\item Graphical User Interface
\end{enumerate}
\answer{D}{Graphical User Interface. GUI provides visual interaction through windows, icons and controls.}
\source{JKSSB official syllabus/question-paper pattern; original synthesis}
\section*{Question 4}
Which is system software?
\begin{enumerate}[label=\Alph*.]
\item Device driver
\item Word document
\item Presentation slide
\item Photograph
\end{enumerate}
\answer{A}{Device driver. A driver helps the OS communicate with hardware.}
\source{JKSSB official syllabus/question-paper pattern; original synthesis}
\section*{Question 5}
Firmware is normally:
\begin{enumerate}[label=\Alph*.]
\item A paper manual
\item Software stored in non-volatile memory for device control
\item A temporary spreadsheet
\item A network cable
\end{enumerate}
\answer{B}{Software stored in non-volatile memory for device control. Firmware provides low-level control of hardware.}
\source{Sanfoundry, Computer Fundamentals chapterwise MCQs; original synthesis}
\section*{Question 6}
Booting means:
\begin{enumerate}[label=\Alph*.]
\item Printing a document
\item Deleting a folder
\item Starting or restarting a computer and loading the OS
\item Compressing a video
\end{enumerate}
\answer{C}{Starting or restarting a computer and loading the OS. Bootstrapping loads the operating system into a usable state.}
\source{JKSSB official syllabus/question-paper pattern; original synthesis}
\section*{Question 7}
Which file operation changes a file’s location but normally keeps its contents?
\begin{enumerate}[label=\Alph*.]
\item Scan
\item Format
\item Compile
\item Move
\end{enumerate}
\answer{D}{Move. Move changes the path; it does not inherently alter file contents.}
\source{JKSSB official syllabus/question-paper pattern; original synthesis}
\section*{Question 8}
A file extension generally indicates:
\begin{enumerate}[label=\Alph*.]
\item The file type or associated format
\item The monitor size
\item The CPU speed
\item The user password
\end{enumerate}
\answer{A}{The file type or associated format. Extensions help the OS and applications identify file formats.}
\source{JKSSB official syllabus/question-paper pattern; original synthesis}
\section*{Question 9}
In Windows, deleted files are commonly first placed in:
\begin{enumerate}[label=\Alph*.]
\item Task Manager
\item Recycle Bin
\item Clipboard
\item Control Panel
\end{enumerate}
\answer{B}{Recycle Bin. The Recycle Bin permits restoration until emptied.}
\source{JKSSB official syllabus/question-paper pattern; original synthesis}
\section*{Question 10}
Which shortcut usually copies selected content in Windows?
\begin{enumerate}[label=\Alph*.]
\item Ctrl+P
\item Ctrl+X
\item Ctrl+C
\item Ctrl+Z
\end{enumerate}
\answer{C}{Ctrl+C. Ctrl+C copies; Ctrl+X cuts; Ctrl+P prints; Ctrl+Z undoes.}
\source{JKSSB official syllabus, MS Office topics; original synthesis}
\section*{Question 11}
Which shortcut usually opens Find and Replace in Office applications?
\begin{enumerate}[label=\Alph*.]
\item Alt+Tab
\item Ctrl+K
\item Ctrl+F5
\item Ctrl+H
\end{enumerate}
\answer{D}{Ctrl+H. Ctrl+H opens Replace in common Office applications.}
\source{JKSSB official syllabus, MS Office topics; original synthesis}
\section*{Question 12}
A real-time operating system is designed to:
\begin{enumerate}[label=\Alph*.]
\item Meet timing constraints predictably
\item Maximize screen resolution
\item Replace all databases
\item Store only images
\end{enumerate}
\answer{A}{Meet timing constraints predictably. RTOS behaviour is evaluated against timing deadlines.}
\source{Sanfoundry, Operating System chapterwise MCQs; original synthesis}
\section*{Question 13}
DMA allows a device to:
\begin{enumerate}[label=\Alph*.]
\item Encrypt every file automatically
\item Transfer data to/from memory with limited CPU intervention
\item Create a user account
\item Render a web page
\end{enumerate}
\answer{B}{Transfer data to/from memory with limited CPU intervention. Direct Memory Access reduces CPU involvement in data transfers.}
\source{JKSSB official syllabus plus current-topic extension; original synthesis}
\section*{Question 14}
Preemptive scheduling allows the OS to:
\begin{enumerate}[label=\Alph*.]
\item Disable the timer
\item Prevent all context switches
\item Interrupt a running process to schedule another
\item Delete the kernel
\end{enumerate}
\answer{C}{Interrupt a running process to schedule another. The scheduler can take the CPU from a running task.}
\source{Sanfoundry, Operating System chapterwise MCQs; original synthesis}
\section*{Question 15}
Starvation occurs when a process:
\begin{enumerate}[label=\Alph*.]
\item Is stored on ROM
\item Finishes immediately
\item Has no instructions
\item Waits indefinitely because others keep receiving resources
\end{enumerate}
\answer{D}{Waits indefinitely because others keep receiving resources. Unfair scheduling or resource allocation can cause indefinite waiting.}
\source{Sanfoundry, Operating System chapterwise MCQs; original synthesis}
\section*{Question 16}
A deadlock recovery method based on taking a resource from a process is:
\begin{enumerate}[label=\Alph*.]
\item Resource preemption
\item Spooling
\item Paging
\item Caching
\end{enumerate}
\answer{A}{Resource preemption. Preemption removes a resource to break a deadlock cycle.}
\source{Sanfoundry, Operating System chapterwise MCQs; original synthesis}
\section*{Question 17}
Linux is best classified as:
\begin{enumerate}[label=\Alph*.]
\item A word processor
\item An open-source operating system kernel/system family
\item A printer interface
\item A database table
\end{enumerate}
\answer{B}{An open-source operating system kernel/system family. Linux is an open-source kernel used in many operating systems.}
\source{Sanfoundry, Operating System chapterwise MCQs; original synthesis}
\section*{Question 18}
If a process is allocated 3 ms of CPU time in each of 4 rounds, its total CPU time is:
\begin{enumerate}[label=\Alph*.]
\item 1 ms
\item 7 ms
\item 12 ms
\item 16 ms
\end{enumerate}
\answer{C}{12 ms. 4 x 3 ms = 12 ms. Numerical/application item.}
\source{Sanfoundry, Operating System chapterwise MCQs; original synthesis}
\section*{Question 19}
Which Windows utility is used to inspect running processes?
\begin{enumerate}[label=\Alph*.]
\item Calculator
\item Paint
\item Notepad
\item Task Manager
\end{enumerate}
\answer{D}{Task Manager. Task Manager shows processes and resource usage.}
\source{JKSSB official syllabus/question-paper pattern; original synthesis}
\section*{Question 20}
An interpreter generally:
\begin{enumerate}[label=\Alph*.]
\item Translates and executes source statements progressively
\item Always creates only hardware
\item Stores files in Recycle Bin
\item Scans paper forms
\end{enumerate}
\answer{A}{Translates and executes source statements progressively. An interpreter executes source code during translation rather than producing a complete object file first.}
\source{Sanfoundry, Computer Fundamentals chapterwise MCQs; original synthesis}

\chapter{MS Word}
\textit{20 questions; numerical/application items are marked in the answer explanation.}
\section*{Question 1}
Which Word feature records insertions, deletions and formatting changes?
\begin{enumerate}[label=\Alph*.]
\item WordArt
\item Mail Merge
\item Track Changes
\item AutoSum
\end{enumerate}
\answer{C}{Track Changes. Track Changes records document revisions.}
\source{JKSSB official syllabus, MS Office topics; original synthesis}
\section*{Question 2}
A page break is used to:
\begin{enumerate}[label=\Alph*.]
\item Insert an email address
\item Change the entire document’s margins only
\item Create a database key
\item Start following content on a new page
\end{enumerate}
\answer{D}{Start following content on a new page. Page breaks control pagination without creating a new section.}
\source{JKSSB official syllabus, MS Office topics; original synthesis}
\section*{Question 3}
A section break can allow different:
\begin{enumerate}[label=\Alph*.]
\item Headers, footers or page layouts within one document
\item CPU registers
\item File systems
\item Network addresses
\end{enumerate}
\answer{A}{Headers, footers or page layouts within one document. Sections permit independent page setup and formatting.}
\source{JKSSB official syllabus, MS Office topics; original synthesis}
\section*{Question 4}
Mail Merge combines a main document with:
\begin{enumerate}[label=\Alph*.]
\item A graphics card
\item A data source and merge fields
\item A printer driver only
\item A web browser cache
\end{enumerate}
\answer{B}{A data source and merge fields. Mail Merge uses records from a data source to produce personalized documents.}
\source{JKSSB official syllabus, MS Office topics; original synthesis}
\section*{Question 5}
Which shortcut opens the Find and Replace dialog in Word?
\begin{enumerate}[label=\Alph*.]
\item Ctrl+M
\item Ctrl+K
\item Ctrl+H
\item Ctrl+W
\end{enumerate}
\answer{C}{Ctrl+H. Ctrl+H invokes Replace in Word.}
\source{JKSSB official syllabus, MS Office topics; original synthesis}
\section*{Question 6}
Which shortcut inserts a hyperlink in Word?
\begin{enumerate}[label=\Alph*.]
\item Ctrl+R
\item Ctrl+H
\item Ctrl+J
\item Ctrl+K
\end{enumerate}
\answer{D}{Ctrl+K. Ctrl+K opens the hyperlink command.}
\source{JKSSB official syllabus, MS Office topics; original synthesis}
\section*{Question 7}
A hanging indent means:
\begin{enumerate}[label=\Alph*.]
\item All lines except the first are indented
\item Only the first line is indented
\item Every line is centred
\item The paragraph is deleted
\end{enumerate}
\answer{A}{All lines except the first are indented. The first line remains at the margin while subsequent lines are indented.}
\source{JKSSB official syllabus, MS Office topics; original synthesis}
\section*{Question 8}
Which view shows a document as it will appear when printed?
\begin{enumerate}[label=\Alph*.]
\item Outline only
\item Print Layout
\item Web Query
\item Formula View
\end{enumerate}
\answer{B}{Print Layout. Print Layout displays page boundaries and print formatting.}
\source{JKSSB official syllabus, MS Office topics; original synthesis}
\section*{Question 9}
The spelling and grammar tool is commonly started with:
\begin{enumerate}[label=\Alph*.]
\item F5
\item F1
\item F7
\item F12
\end{enumerate}
\answer{C}{F7. F7 is the standard Word spelling/grammar shortcut.}
\source{JKSSB official syllabus, MS Office topics; original synthesis}
\section*{Question 10}
A header appears:
\begin{enumerate}[label=\Alph*.]
\item Inside a network packet
\item Only inside a spreadsheet cell
\item At the bottom of a printer
\item At the top area of document pages
\end{enumerate}
\answer{D}{At the top area of document pages. Headers are repeating top-of-page content.}
\source{JKSSB official syllabus, MS Office topics; original synthesis}
\section*{Question 11}
Styles in Word are useful for:
\begin{enumerate}[label=\Alph*.]
\item Applying consistent formatting and navigation structure
\item Increasing CPU clock speed
\item Encrypting Wi-Fi
\item Converting HDMI to VGA
\end{enumerate}
\answer{A}{Applying consistent formatting and navigation structure. Styles standardize headings and other formatting.}
\source{JKSSB official syllabus, MS Office topics; original synthesis}
\section*{Question 12}
The Navigation Pane can help users:
\begin{enumerate}[label=\Alph*.]
\item Format a hard disk
\item Find headings, pages or search results
\item Set an IP address
\item Install RAM
\end{enumerate}
\answer{B}{Find headings, pages or search results. It provides document navigation and search.}
\source{JKSSB official syllabus, MS Office topics; original synthesis}
\section*{Question 13}
Comments are primarily used to:
\begin{enumerate}[label=\Alph*.]
\item Create a pivot table
\item Replace the operating system
\item Add review notes without changing main text
\item Send an SMTP message
\end{enumerate}
\answer{C}{Add review notes without changing main text. Comments support collaborative review.}
\source{JKSSB official syllabus, MS Office topics; original synthesis}
\section*{Question 14}
PDF export from Word generally creates:
\begin{enumerate}[label=\Alph*.]
\item A network protocol
\item A CPU instruction
\item A database relationship
\item A fixed-layout sharing format
\end{enumerate}
\answer{D}{A fixed-layout sharing format. PDF preserves layout for distribution and printing.}
\source{JKSSB official syllabus, MS Office topics; original synthesis}
\section*{Question 15}
Which alignment makes both left and right paragraph edges appear even?
\begin{enumerate}[label=\Alph*.]
\item Justify
\item Left
\item Right
\item Centre
\end{enumerate}
\answer{A}{Justify. Justify adjusts spacing so both edges align.}
\source{JKSSB official syllabus, MS Office topics; original synthesis}
\section*{Question 16}
Bullets are used primarily to:
\begin{enumerate}[label=\Alph*.]
\item Calculate averages
\item Create an unordered list
\item Manage processes
\item Resolve DNS
\end{enumerate}
\answer{B}{Create an unordered list. Bullets visually mark list items without ranking them numerically.}
\source{JKSSB official syllabus, MS Office topics; original synthesis}
\section*{Question 17}
A document has 4 pages with 2 header fields per page. How many header instances are displayed if each page uses the same header?
\begin{enumerate}[label=\Alph*.]
\item 6
\item 2
\item 4
\item 8
\end{enumerate}
\answer{C}{4. The same header is displayed once on each of 4 pages. Numerical/application item.}
\source{JKSSB official syllabus, MS Office topics; original synthesis}
\section*{Question 18}
A bookmark in Word identifies:
\begin{enumerate}[label=\Alph*.]
\item A network switch
\item A printer cartridge
\item An Excel formula
\item A named location for quick navigation
\end{enumerate}
\answer{D}{A named location for quick navigation. Bookmarks mark locations or text for links and navigation.}
\source{JKSSB official syllabus, MS Office topics; original synthesis}
\section*{Question 19}
Which is not normally a Word document extension?
\begin{enumerate}[label=\Alph*.]
\item .docx
\item .docm
\item .dotx
\item .xlsx
\end{enumerate}
\answer{A}{.docx. .xlsx is the standard Excel workbook extension.}
\source{JKSSB official syllabus, MS Office topics; original synthesis}
\section*{Question 20}
Track Changes metadata may show:
\begin{enumerate}[label=\Alph*.]
\item The CPU temperature only
\item Who made a revision and when
\item The Wi-Fi channel only
\item The printer toner level
\end{enumerate}
\answer{B}{Who made a revision and when. Revision metadata can identify the reviewer and change time.}
\source{JKSSB official syllabus, MS Office topics; original synthesis}

\chapter{MS Excel}
\textit{20 questions; numerical/application items are marked in the answer explanation.}
\section*{Question 1}
The intersection of a row and column in Excel is a:
\begin{enumerate}[label=\Alph*.]
\item Worksheet tab only
\item Workbook
\item Chart
\item Cell
\end{enumerate}
\answer{D}{Cell. A cell is identified by a column letter and row number.}
\source{JKSSB official syllabus, MS Office topics; original synthesis}
\section*{Question 2}
A group of worksheets saved together is a:
\begin{enumerate}[label=\Alph*.]
\item Workbook
\item Range
\item Formula bar
\item Name Box
\end{enumerate}
\answer{A}{Workbook. An Excel file is a workbook containing worksheets.}
\source{JKSSB official syllabus, MS Office topics; original synthesis}
\section*{Question 3}
Which reference remains fixed when copied?
\begin{enumerate}[label=\Alph*.]
\item A1
\item \$A\$1
\item A\$1 only
\item \$A1 only
\end{enumerate}
\answer{B}{\$A\$1. Dollar signs before both column and row make an absolute reference.}
\source{JKSSB official syllabus, MS Office topics; original synthesis}
\section*{Question 4}
Which function calculates the arithmetic mean?
\begin{enumerate}[label=\Alph*.]
\item MAX
\item COUNT
\item AVERAGE
\item IF
\end{enumerate}
\answer{C}{AVERAGE. AVERAGE returns the mean of numeric arguments.}
\source{JKSSB official syllabus, MS Office topics; original synthesis}
\section*{Question 5}
COUNT counts:
\begin{enumerate}[label=\Alph*.]
\item Only worksheets
\item All non-empty cells including text
\item Only blank cells
\item Numeric cells
\end{enumerate}
\answer{D}{Numeric cells. COUNT counts cells containing numbers.}
\source{JKSSB official syllabus, MS Office topics; original synthesis}
\section*{Question 6}
COUNTA counts:
\begin{enumerate}[label=\Alph*.]
\item Non-empty cells
\item Only numeric cells
\item Only formulas returning zero
\item Only hidden rows
\end{enumerate}
\answer{A}{Non-empty cells. COUNTA counts cells that are not empty.}
\source{JKSSB official syllabus, MS Office topics; original synthesis}
\section*{Question 7}
Which function returns the largest value in a range?
\begin{enumerate}[label=\Alph*.]
\item MIN
\item MAX
\item SUM
\item COUNTIF
\end{enumerate}
\answer{B}{MAX. MAX returns the largest numeric value.}
\source{JKSSB official syllabus, MS Office topics; original synthesis}
\section*{Question 8}
The IF function is used for:
\begin{enumerate}[label=\Alph*.]
\item Sending email
\item Sorting a hard disk
\item Logical conditional evaluation
\item Changing screen resolution
\end{enumerate}
\answer{C}{Logical conditional evaluation. IF returns one result when a condition is true and another when false.}
\source{JKSSB official syllabus, MS Office topics; original synthesis}
\section*{Question 9}
A formula in Excel normally begins with:
\begin{enumerate}[label=\Alph*.]
\item @
\item +
\item \#
\item =
\end{enumerate}
\answer{D}{=. The equals sign starts a formula.}
\source{JKSSB official syllabus, MS Office topics; original synthesis}
\section*{Question 10}
Which tool displays only rows meeting criteria?
\begin{enumerate}[label=\Alph*.]
\item Filter
\item Merge Cells
\item Freeze Panes
\item Page Break
\end{enumerate}
\answer{A}{Filter. Filtering hides rows that do not meet the criteria.}
\source{JKSSB official syllabus, MS Office topics; original synthesis}
\section*{Question 11}
A PivotTable is primarily used to:
\begin{enumerate}[label=\Alph*.]
\item Edit video frames
\item Summarize and analyze tabular data
\item Install drivers
\item Create an operating system
\end{enumerate}
\answer{B}{Summarize and analyze tabular data. PivotTables aggregate and reorganize data.}
\source{JKSSB official syllabus, MS Office topics; original synthesis}
\section*{Question 12}
A chart is used to:
\begin{enumerate}[label=\Alph*.]
\item Manage CPU processes
\item Compile source code
\item Visualize data
\item Resolve URLs
\end{enumerate}
\answer{C}{Visualize data. Charts present data graphically.}
\source{JKSSB official syllabus, MS Office topics; original synthesis}
\section*{Question 13}
The formula =SUM(A1:A3) adds:
\begin{enumerate}[label=\Alph*.]
\item All worksheets
\item Only A1 and A3
\item The cell addresses as text
\item The values in A1 through A3
\end{enumerate}
\answer{D}{The values in A1 through A3. A1:A3 is a contiguous range of three cells.}
\source{JKSSB official syllabus, MS Office topics; original synthesis}
\section*{Question 14}
If A1=10 and A2=15, what is =SUM(A1:A2)?
\begin{enumerate}[label=\Alph*.]
\item 25
\item 5
\item 150
\item 1015
\end{enumerate}
\answer{A}{25. 10 + 15 = 25. Numerical/application item.}
\source{JKSSB official syllabus, MS Office topics; original synthesis}
\section*{Question 15}
Which error commonly indicates division by zero?
\begin{enumerate}[label=\Alph*.]
\item \#NAME?
\item \#DIV/0!
\item \#REF!
\item \#VALUE!
\end{enumerate}
\answer{B}{\#DIV/0!. \#DIV/0! occurs when a formula divides by zero.}
\source{JKSSB official syllabus, MS Office topics; original synthesis}
\section*{Question 16}
A mixed reference is:
\begin{enumerate}[label=\Alph*.]
\item A1 only
\item \$A\$1 only
\item A\$1 or \$A1
\item A:A only
\end{enumerate}
\answer{C}{A\$1 or \$A1. One component is fixed and the other is relative.}
\source{JKSSB official syllabus, MS Office topics; original synthesis}
\section*{Question 17}
Sorting changes:
\begin{enumerate}[label=\Alph*.]
\item The worksheet name necessarily
\item The CPU architecture
\item The file extension automatically
\item The order of records according to selected values
\end{enumerate}
\answer{D}{The order of records according to selected values. Sorting rearranges data order.}
\source{JKSSB official syllabus, MS Office topics; original synthesis}
\section*{Question 18}
COUNTIFS supports:
\begin{enumerate}[label=\Alph*.]
\item Counting using multiple criteria
\item Only formatting cells
\item Creating slides
\item Sending attachments
\end{enumerate}
\answer{A}{Counting using multiple criteria. COUNTIFS counts records satisfying multiple conditions.}
\source{JKSSB official syllabus plus current-topic extension; original synthesis}
\section*{Question 19}
SUMPRODUCT generally calculates:
\begin{enumerate}[label=\Alph*.]
\item A text spelling score
\item The sum of products of corresponding array elements
\item An IP address
\item A slide transition
\end{enumerate}
\answer{B}{The sum of products of corresponding array elements. SUMPRODUCT multiplies corresponding elements and sums the products.}
\source{JKSSB official syllabus plus current-topic extension; original synthesis}
\section*{Question 20}
If a worksheet has 5 rows and 4 columns of numeric values, how many numeric cells are present?
\begin{enumerate}[label=\Alph*.]
\item 16
\item 9
\item 20
\item 25
\end{enumerate}
\answer{C}{20. 5 x 4 = 20. Numerical/application item.}
\source{JKSSB official syllabus, MS Office topics; original synthesis}

\chapter{MS Access and Databases}
\textit{20 questions; numerical/application items are marked in the answer explanation.}
\section*{Question 1}
A database is:
\begin{enumerate}[label=\Alph*.]
\item An organized collection of related data
\item A display cable
\item A CPU register
\item A printer driver
\end{enumerate}
\answer{A}{An organized collection of related data. Databases store related data for access and management.}
\source{JKSSB official syllabus/question-paper pattern; original synthesis}
\section*{Question 2}
A DBMS is software used to:
\begin{enumerate}[label=\Alph*.]
\item Control monitor brightness only
\item Create, store, retrieve and manage databases
\item Replace a keyboard
\item Compress images only
\end{enumerate}
\answer{B}{Create, store, retrieve and manage databases. DBMS software manages database structures and operations.}
\source{JKSSB official syllabus/question-paper pattern; original synthesis}
\section*{Question 3}
In a relational table, a row is commonly called a:
\begin{enumerate}[label=\Alph*.]
\item Query
\item Field
\item Record
\item Form
\end{enumerate}
\answer{C}{Record. A record represents one row/entity instance.}
\source{JKSSB official syllabus/question-paper pattern; original synthesis}
\section*{Question 4}
In a relational table, a column is commonly called a:
\begin{enumerate}[label=\Alph*.]
\item Relationship
\item Record
\item Report
\item Field
\end{enumerate}
\answer{D}{Field. A field stores one attribute across records.}
\source{JKSSB official syllabus/question-paper pattern; original synthesis}
\section*{Question 5}
A primary key must:
\begin{enumerate}[label=\Alph*.]
\item Uniquely identify each record
\item Contain duplicate values necessarily
\item Always be a password
\item Only contain images
\end{enumerate}
\answer{A}{Uniquely identify each record. A primary key provides unique record identification.}
\source{JKSSB official syllabus/question-paper pattern; original synthesis}
\section*{Question 6}
A foreign key is used mainly to:
\begin{enumerate}[label=\Alph*.]
\item Format a document
\item Link a record to a key in another table
\item Run a presentation
\item Control a monitor
\end{enumerate}
\answer{B}{Link a record to a key in another table. Foreign keys represent relationships between tables.}
\source{JKSSB official syllabus/question-paper pattern; original synthesis}
\section*{Question 7}
A key made from more than one attribute is a:
\begin{enumerate}[label=\Alph*.]
\item Simple field
\item Foreign monitor
\item Composite key
\item Report key only
\end{enumerate}
\answer{C}{Composite key. Composite keys combine multiple columns.}
\source{JKSSB official syllabus/question-paper pattern; original synthesis}
\section*{Question 8}
A one-to-many relationship means:
\begin{enumerate}[label=\Alph*.]
\item Only many tables exist
\item Every record relates to exactly one record only
\item No records are related
\item One record can relate to many records in another table
\end{enumerate}
\answer{D}{One record can relate to many records in another table. For example, one customer may have many orders.}
\source{JKSSB official syllabus/question-paper pattern; original synthesis}
\section*{Question 9}
A junction table is commonly used to implement:
\begin{enumerate}[label=\Alph*.]
\item Many-to-many relationships
\item A single text field
\item A printer queue
\item A slide layout
\end{enumerate}
\answer{A}{Many-to-many relationships. The junction table stores pairs of related keys.}
\source{JKSSB official syllabus/question-paper pattern; original synthesis}
\section*{Question 10}
Referential integrity helps ensure that:
\begin{enumerate}[label=\Alph*.]
\item Every field is encrypted
\item Foreign-key references remain valid
\item All tables have pictures
\item The monitor never flickers
\end{enumerate}
\answer{B}{Foreign-key references remain valid. It prevents invalid or orphaned relationships.}
\source{JKSSB official syllabus/question-paper pattern; original synthesis}
\section*{Question 11}
A query is used to:
\begin{enumerate}[label=\Alph*.]
\item Print a keyboard
\item Draw a hardware circuit
\item Retrieve or manipulate selected data
\item Boot the system
\end{enumerate}
\answer{C}{Retrieve or manipulate selected data. Queries select, filter, calculate or modify database data.}
\source{JKSSB official syllabus/question-paper pattern; original synthesis}
\section*{Question 12}
A form is mainly used to:
\begin{enumerate}[label=\Alph*.]
\item Compile programs
\item Store the operating system kernel
\item Route packets
\item Enter and view data conveniently
\end{enumerate}
\answer{D}{Enter and view data conveniently. Forms provide a user interface for data entry and viewing.}
\source{JKSSB official syllabus/question-paper pattern; original synthesis}
\section*{Question 13}
A report is mainly used to:
\begin{enumerate}[label=\Alph*.]
\item Present formatted database information
\item Change CPU voltage
\item Recognize speech
\item Assign an IP address
\end{enumerate}
\answer{A}{Present formatted database information. Reports format data for viewing or printing.}
\source{JKSSB official syllabus/question-paper pattern; original synthesis}
\section*{Question 14}
An index is used mainly to:
\begin{enumerate}[label=\Alph*.]
\item Increase monitor size
\item Speed up data retrieval
\item Replace all keys
\item Create a slide transition
\end{enumerate}
\answer{B}{Speed up data retrieval. Indexes improve lookup performance at a storage cost.}
\source{JKSSB official syllabus/question-paper pattern; original synthesis}
\section*{Question 15}
Which is a common MS Access database extension?
\begin{enumerate}[label=\Alph*.]
\item .pptx
\item .xlsx
\item .accdb
\item .mp3
\end{enumerate}
\answer{C}{.accdb. .accdb is the modern Access database format.}
\source{JKSSB official syllabus/question-paper pattern; original synthesis}
\section*{Question 16}
If a table has 6 records and each record has 4 fields, how many field values are stored, ignoring nulls?
\begin{enumerate}[label=\Alph*.]
\item 64
\item 10
\item 2
\item 24
\end{enumerate}
\answer{D}{24. 6 x 4 = 24 field positions. Numerical/application item.}
\source{JKSSB official syllabus/question-paper pattern; original synthesis}
\section*{Question 17}
A data type defines:
\begin{enumerate}[label=\Alph*.]
\item The kind of value a field can store
\item The screen refresh rate
\item The network topology
\item The printer speed
\end{enumerate}
\answer{A}{The kind of value a field can store. Data types constrain and describe field values.}
\source{JKSSB official syllabus/question-paper pattern; original synthesis}
\section*{Question 18}
Validation rules are used to:
\begin{enumerate}[label=\Alph*.]
\item Create a CPU core
\item Restrict unacceptable data entry
\item Open a web browser
\item Convert a PDF to audio
\end{enumerate}
\answer{B}{Restrict unacceptable data entry. Validation protects data quality at entry.}
\source{JKSSB official syllabus/question-paper pattern; original synthesis}
\section*{Question 19}
Which sequence reflects a common Access workflow?
\begin{enumerate}[label=\Alph*.]
\item Form–Printer–Kernel–Query
\item Report–CPU–Table–RAM
\item Table–Query–Form–Report
\item Query–Monitor–Table–BIOS
\end{enumerate}
\answer{C}{Table–Query–Form–Report. Tables store, queries select, forms interact and reports present.}
\source{JKSSB official syllabus/question-paper pattern; original synthesis}
\section*{Question 20}
A many-to-many relationship between students and courses normally requires:
\begin{enumerate}[label=\Alph*.]
\item A cache register
\item Only one text box
\item A printer queue
\item A junction/enrolment table
\end{enumerate}
\answer{D}{A junction/enrolment table. The junction table stores student-course pairs.}
\source{JKSSB official syllabus/question-paper pattern; original synthesis}

\chapter{PowerPoint, PDF and Multimedia}
\textit{20 questions; numerical/application items are marked in the answer explanation.}
\section*{Question 1}
A PowerPoint slide layout defines:
\begin{enumerate}[label=\Alph*.]
\item The CPU instruction set
\item Positions of placeholders and content areas
\item A database index
\item An email protocol
\end{enumerate}
\answer{B}{Positions of placeholders and content areas. Layouts provide predefined slide structures.}
\source{JKSSB official syllabus, MS Office topics; original synthesis}
\section*{Question 2}
Slide Master controls:
\begin{enumerate}[label=\Alph*.]
\item Email delivery
\item Hard-disk partitions
\item Consistent design and formatting across slides
\item RAM voltage
\end{enumerate}
\answer{C}{Consistent design and formatting across slides. The Slide Master applies common elements and formatting.}
\source{JKSSB official syllabus, MS Office topics; original synthesis}
\section*{Question 3}
An animation is applied primarily to:
\begin{enumerate}[label=\Alph*.]
\item An IP address
\item The entire computer boot process
\item A database key
\item Objects on a slide
\end{enumerate}
\answer{D}{Objects on a slide. Animations affect text or objects on a slide.}
\source{JKSSB official syllabus, MS Office topics; original synthesis}
\section*{Question 4}
A transition occurs:
\begin{enumerate}[label=\Alph*.]
\item When moving from one slide to another
\item Only while editing an Excel cell
\item During CPU fetch
\item When a file is renamed
\end{enumerate}
\answer{A}{When moving from one slide to another. Transitions control slide-to-slide movement.}
\source{JKSSB official syllabus, MS Office topics; original synthesis}
\section*{Question 5}
F5 in PowerPoint normally starts the slide show:
\begin{enumerate}[label=\Alph*.]
\item From the current slide only
\item From the beginning
\item In the File Explorer
\item In print preview only
\end{enumerate}
\answer{B}{From the beginning. F5 starts from the first slide; Shift+F5 starts from the current slide.}
\source{JKSSB official syllabus, MS Office topics; original synthesis}
\section*{Question 6}
Shift+F5 in PowerPoint normally starts the show:
\begin{enumerate}[label=\Alph*.]
\item In outline view only
\item From the first slide
\item From the current slide
\item As a database query
\end{enumerate}
\answer{C}{From the current slide. Shift+F5 starts from the current slide.}
\source{JKSSB official syllabus, MS Office topics; original synthesis}
\section*{Question 7}
PDF is primarily designed as a:
\begin{enumerate}[label=\Alph*.]
\item Spreadsheet formula
\item CPU scheduling algorithm
\item Network cable type
\item Fixed-layout document format
\end{enumerate}
\answer{D}{Fixed-layout document format. PDF aims to preserve document appearance across systems.}
\source{JKSSB official syllabus, MS Office topics; original synthesis}
\section*{Question 8}
A scanned PDF may contain:
\begin{enumerate}[label=\Alph*.]
\item Only images without selectable text
\item Only formulas
\item Only audio
\item A live database connection necessarily
\end{enumerate}
\answer{A}{Only images without selectable text. Scanned pages are image data unless OCR is applied.}
\source{JKSSB official syllabus, MS Office topics; original synthesis}
\section*{Question 9}
OCR can make a scanned document:
\begin{enumerate}[label=\Alph*.]
\item A network switch
\item Searchable or editable as recognized text
\item An operating system
\item A spreadsheet chart only
\end{enumerate}
\answer{B}{Searchable or editable as recognized text. OCR converts image characters into machine-readable text.}
\source{JKSSB official syllabus, MS Office topics; original synthesis}
\section*{Question 10}
PDF/A is intended primarily for:
\begin{enumerate}[label=\Alph*.]
\item CPU cache management
\item Real-time gaming
\item Long-term archival preservation
\item Email spam filtering
\end{enumerate}
\answer{C}{Long-term archival preservation. PDF/A constrains features to improve archival reliability.}
\source{JKSSB official syllabus, MS Office topics; original synthesis}
\section*{Question 11}
A digital signature helps provide:
\begin{enumerate}[label=\Alph*.]
\item A faster keyboard
\item A larger monitor
\item More RAM
\item Authenticity and integrity evidence
\end{enumerate}
\answer{D}{Authenticity and integrity evidence. Signatures can verify signer identity and detect changes.}
\source{JKSSB official syllabus plus current-topic extension; original synthesis}
\section*{Question 12}
MP3 is primarily an:
\begin{enumerate}[label=\Alph*.]
\item Audio format
\item Image format
\item Spreadsheet format
\item Database format
\end{enumerate}
\answer{A}{Audio format. MP3 commonly stores compressed audio.}
\source{JKSSB official syllabus, MS Office topics; original synthesis}
\section*{Question 13}
MP4 is commonly associated with:
\begin{enumerate}[label=\Alph*.]
\item Font files only
\item Video container/media
\item Database tables
\item Executable code only
\end{enumerate}
\answer{B}{Video container/media. MP4 is widely used for digital video and audio.}
\source{JKSSB official syllabus, MS Office topics; original synthesis}
\section*{Question 14}
PNG is primarily an:
\begin{enumerate}[label=\Alph*.]
\item Operating system
\item Email protocol
\item Image format
\item Printer language
\end{enumerate}
\answer{C}{Image format. PNG is a raster image format.}
\source{JKSSB official syllabus, MS Office topics; original synthesis}
\section*{Question 15}
WYSIWYG means:
\begin{enumerate}[label=\Alph*.]
\item Write Your Source in Word Globally
\item Where Your System Is Working Quickly
\item Web Yield System in Web Graphics
\item What You See Is What You Get
\end{enumerate}
\answer{D}{What You See Is What You Get. The displayed design approximates the output.}
\source{JKSSB official syllabus, MS Office topics; original synthesis}
\section*{Question 16}
Multimedia digitization commonly follows:
\begin{enumerate}[label=\Alph*.]
\item Sampling–quantization–encoding–compression
\item Printing–scanning–booting–routing
\item Sorting–filtering–joining–pivoting
\item Fetching–decoding–compiling–paging
\end{enumerate}
\answer{A}{Sampling–quantization–encoding–compression. These stages convert and prepare analogue media for digital use.}
\source{JKSSB official syllabus plus current-topic extension; original synthesis}
\section*{Question 17}
A presentation has 8 slides and each slide contains 3 main objects. How many objects are there in total?
\begin{enumerate}[label=\Alph*.]
\item 11
\item 24
\item 16
\item 32
\end{enumerate}
\answer{B}{24. 8 x 3 = 24. Numerical/application item.}
\source{JKSSB official syllabus, MS Office topics; original synthesis}
\section*{Question 18}
Speaker notes are used to:
\begin{enumerate}[label=\Alph*.]
\item Create a primary key
\item Change the CPU clock
\item Store presenter guidance associated with slides
\item Set a DNS record
\end{enumerate}
\answer{C}{Store presenter guidance associated with slides. Notes help the presenter and need not appear on the slide.}
\source{JKSSB official syllabus, MS Office topics; original synthesis}
\section*{Question 19}
A handout is designed primarily for:
\begin{enumerate}[label=\Alph*.]
\item Routing packets
\item Storing RAM
\item Editing database rows
\item Printing multiple slide thumbnails for an audience
\end{enumerate}
\answer{D}{Printing multiple slide thumbnails for an audience. Handouts distribute slide content in printed form.}
\source{JKSSB official syllabus, MS Office topics; original synthesis}
\section*{Question 20}
Redaction in a PDF is used to:
\begin{enumerate}[label=\Alph*.]
\item Permanently remove sensitive visible content when properly applied
\item Add animation
\item Increase audio sampling rate
\item Create a spreadsheet formula
\end{enumerate}
\answer{A}{Permanently remove sensitive visible content when properly applied. Proper redaction removes or obscures sensitive content rather than merely covering it.}
\source{JKSSB official syllabus plus current-topic extension; original synthesis}

\chapter{Internet, Web, Email, E-Banking and Cloud}
\textit{20 questions; numerical/application items are marked in the answer explanation.}
\section*{Question 1}
The Internet is best described as:
\begin{enumerate}[label=\Alph*.]
\item A word processor
\item A single website
\item A global network of interconnected networks
\item A printer driver
\end{enumerate}
\answer{C}{A global network of interconnected networks. The Internet links many independent networks.}
\source{Sanfoundry, Computer Network chapterwise MCQs; original synthesis}
\section*{Question 2}
The World Wide Web is:
\begin{enumerate}[label=\Alph*.]
\item A spreadsheet function
\item Identical to the CPU
\item A type of RAM
\item A service using linked resources over the Internet
\end{enumerate}
\answer{D}{A service using linked resources over the Internet. The Web is one Internet service based on web resources and protocols.}
\source{Sanfoundry, Computer Network chapterwise MCQs; original synthesis}
\section*{Question 3}
A browser is used to:
\begin{enumerate}[label=\Alph*.]
\item Retrieve and display web resources
\item Compile a kernel necessarily
\item Print bank cheques
\item Manage RAM chips
\end{enumerate}
\answer{A}{Retrieve and display web resources. Browsers request and render web pages and other resources.}
\source{Sanfoundry, Computer Network chapterwise MCQs; original synthesis}
\section*{Question 4}
A search engine primarily:
\begin{enumerate}[label=\Alph*.]
\item Replaces the ISP
\item Crawls, indexes and ranks web content
\item Converts RAM to ROM
\item Formats a hard disk
\end{enumerate}
\answer{B}{Crawls, indexes and ranks web content. Search systems organize indexed content for retrieval.}
\source{Sanfoundry, Computer Network chapterwise MCQs; original synthesis}
\section*{Question 5}
DNS maps domain names to:
\begin{enumerate}[label=\Alph*.]
\item Printer pixels
\item Keyboard scan codes
\item IP addresses or other DNS records
\item Excel formulas
\end{enumerate}
\answer{C}{IP addresses or other DNS records. DNS resolves human-readable names through records.}
\source{Sanfoundry, Computer Network chapterwise MCQs; original synthesis}
\section*{Question 6}
HTTPS adds security to HTTP using:
\begin{enumerate}[label=\Alph*.]
\item BIOS
\item FTP
\item SMTP
\item TLS
\end{enumerate}
\answer{D}{TLS. HTTPS is HTTP protected by TLS in common usage.}
\source{Sanfoundry, Computer Network chapterwise MCQs; original synthesis}
\section*{Question 7}
Which URL component identifies the communication scheme?
\begin{enumerate}[label=\Alph*.]
\item https
\item example.com
\item /path
\item ?q=term
\end{enumerate}
\answer{A}{https. The scheme appears before the colon in a URL.}
\source{Sanfoundry, Computer Network chapterwise MCQs; original synthesis}
\section*{Question 8}
Uploading means:
\begin{enumerate}[label=\Alph*.]
\item Receiving a file from a server
\item Sending data from a local device to a remote system
\item Deleting a browser tab
\item Encrypting a CPU register
\end{enumerate}
\answer{B}{Sending data from a local device to a remote system. Upload direction is local to remote.}
\source{Sanfoundry, Computer Network chapterwise MCQs; original synthesis}
\section*{Question 9}
In email, Bcc recipients are:
\begin{enumerate}[label=\Alph*.]
\item The only recipients who can reply
\item Always unable to receive attachments
\item Hidden from other recipients in the message header
\item Network administrators only
\end{enumerate}
\answer{C}{Hidden from other recipients in the message header. Bcc conceals those addresses from other recipients.}
\source{JKSSB official syllabus/question-paper pattern; original synthesis}
\section*{Question 10}
SMTP is used primarily for:
\begin{enumerate}[label=\Alph*.]
\item Assigning IP addresses
\item Receiving web pages
\item Resolving domain names
\item Sending email
\end{enumerate}
\answer{D}{Sending email. SMTP transfers outgoing email between clients and servers/servers.}
\source{Sanfoundry, Computer Network chapterwise MCQs; original synthesis}
\section*{Question 11}
IMAP is useful because it:
\begin{enumerate}[label=\Alph*.]
\item Keeps mail synchronized on the server across clients
\item Only prints email
\item Replaces DNS
\item Encrypts every attachment automatically
\end{enumerate}
\answer{A}{Keeps mail synchronized on the server across clients. IMAP supports server-side mailbox synchronization.}
\source{Sanfoundry, Computer Network chapterwise MCQs; original synthesis}
\section*{Question 12}
POP3 traditionally emphasizes:
\begin{enumerate}[label=\Alph*.]
\item Routing packets between autonomous systems
\item Downloading mail to a client
\item Rendering HTML
\item Managing spreadsheet cells
\end{enumerate}
\answer{B}{Downloading mail to a client. POP3 is a mail retrieval protocol traditionally oriented to download.}
\source{Sanfoundry, Computer Network chapterwise MCQs; original synthesis}
\section*{Question 13}
Phishing is:
\begin{enumerate}[label=\Alph*.]
\item A type of monitor
\item A legitimate backup method
\item A deceptive attempt to obtain information
\item A spreadsheet chart
\end{enumerate}
\answer{C}{A deceptive attempt to obtain information. Phishing uses deception, often through messages or fake sites.}
\source{JKSSB official syllabus plus current-topic extension; original synthesis}
\section*{Question 14}
An OTP should be:
\begin{enumerate}[label=\Alph*.]
\item Printed on every email
\item Posted publicly
\item Used as a username
\item Kept secret and never shared with callers
\end{enumerate}
\answer{D}{Kept secret and never shared with callers. One-time passwords are authentication secrets.}
\source{JKSSB official syllabus plus current-topic extension; original synthesis}
\section*{Question 15}
UPI PIN is used to:
\begin{enumerate}[label=\Alph*.]
\item Authorize UPI transactions
\item Identify a printer model
\item Open a PDF only
\item Change screen resolution
\end{enumerate}
\answer{A}{Authorize UPI transactions. The UPI PIN authorizes payment transactions.}
\source{JKSSB official syllabus plus current-topic extension; original synthesis}
\section*{Question 16}
IaaS provides primarily:
\begin{enumerate}[label=\Alph*.]
\item Only word-processing templates
\item Virtualized compute, storage and networking infrastructure
\item A physical keyboard to each user
\item A printed certificate
\end{enumerate}
\answer{B}{Virtualized compute, storage and networking infrastructure. Infrastructure as a Service supplies foundational cloud resources.}
\source{JKSSB official syllabus plus current-topic extension; original synthesis}
\section*{Question 17}
TCP is designed to provide:
\begin{enumerate}[label=\Alph*.]
\item No sequencing
\item Only broadcast video
\item Reliable, ordered delivery
\item A display connector
\end{enumerate}
\answer{C}{Reliable, ordered delivery. TCP uses acknowledgements and sequencing for reliable delivery.}
\source{Sanfoundry, Computer Network chapterwise MCQs; original synthesis}
\section*{Question 18}
If an email message is 2 MB and an attachment is 3 MB, ignoring overhead, the total is:
\begin{enumerate}[label=\Alph*.]
\item 2.5 MB
\item 1 MB
\item 6 MB
\item 5 MB
\end{enumerate}
\answer{D}{5 MB. 2 + 3 = 5 MB. Numerical/application item.}
\source{JKSSB official syllabus/question-paper pattern; original synthesis}
\section*{Question 19}
Cookies are commonly used by websites to:
\begin{enumerate}[label=\Alph*.]
\item Store small pieces of client-side/session-related data
\item Increase CPU cores
\item Replace antivirus software
\item Convert PDF to DOCX without software
\end{enumerate}
\answer{A}{Store small pieces of client-side/session-related data. Cookies can preserve preferences or session identifiers.}
\source{Sanfoundry, Computer Network chapterwise MCQs; original synthesis}
\section*{Question 20}
QUIC is commonly associated with:
\begin{enumerate}[label=\Alph*.]
\item Email over POP3 only
\item HTTP/3 transport over UDP
\item The BIOS boot sequence
\item A printer connector
\end{enumerate}
\answer{B}{HTTP/3 transport over UDP. HTTP/3 commonly uses QUIC, which runs over UDP.}
\source{JKSSB official syllabus plus current-topic extension; original synthesis}

\chapter{Networking, Cybersecurity, Open Source and Governance}
\textit{20 questions; numerical/application items are marked in the answer explanation.}
\section*{Question 1}
A LAN usually covers:
\begin{enumerate}[label=\Alph*.]
\item A satellite orbit only
\item The entire world necessarily
\item Only one CPU register
\item A limited local area such as a building
\end{enumerate}
\answer{D}{A limited local area such as a building. LANs operate over relatively small geographic areas.}
\source{Sanfoundry, Computer Network chapterwise MCQs; original synthesis}
\section*{Question 2}
A router primarily:
\begin{enumerate}[label=\Alph*.]
\item Forwards packets between networks
\item Connects only keyboard keys
\item Prints documents
\item Stores Word styles
\end{enumerate}
\answer{A}{Forwards packets between networks. Routers make forwarding decisions between networks.}
\source{Sanfoundry, Computer Network chapterwise MCQs; original synthesis}
\section*{Question 3}
A switch commonly forwards frames using:
\begin{enumerate}[label=\Alph*.]
\item Only domain names
\item MAC addresses
\item Spreadsheet formulas
\item File extensions
\end{enumerate}
\answer{B}{MAC addresses. Ethernet switches learn and use MAC addresses.}
\source{Sanfoundry, Computer Network chapterwise MCQs; original synthesis}
\section*{Question 4}
A hub generally:
\begin{enumerate}[label=\Alph*.]
\item Encrypts files
\item Makes routing decisions using BGP
\item Repeats incoming signals to all ports
\item Stores DNS zones
\end{enumerate}
\answer{C}{Repeats incoming signals to all ports. A hub is a simple multiport repeater.}
\source{Sanfoundry, Computer Network chapterwise MCQs; original synthesis}
\section*{Question 5}
DHCP commonly provides:
\begin{enumerate}[label=\Alph*.]
\item Keyboard shortcuts
\item Email encryption
\item PDF redaction
\item Automatic IP configuration
\end{enumerate}
\answer{D}{Automatic IP configuration. DHCP leases network configuration to clients.}
\source{Sanfoundry, Computer Network chapterwise MCQs; original synthesis}
\section*{Question 6}
A firewall is used to:
\begin{enumerate}[label=\Alph*.]
\item Control traffic according to security rules
\item Increase RAM capacity
\item Create a slide master
\item Convert analog sound
\end{enumerate}
\answer{A}{Control traffic according to security rules. Firewalls enforce traffic policies at boundaries or hosts.}
\source{JKSSB official syllabus plus current-topic extension; original synthesis}
\section*{Question 7}
A virus generally requires:
\begin{enumerate}[label=\Alph*.]
\item No host ever
\item A host file or program to replicate
\item Only a printer
\item A DNS record
\end{enumerate}
\answer{B}{A host file or program to replicate. Traditional viruses attach to or modify host programs/files.}
\source{JKSSB official syllabus plus current-topic extension; original synthesis}
\section*{Question 8}
A worm differs from a traditional virus because it can:
\begin{enumerate}[label=\Alph*.]
\item Never use networks
\item Only infect paper
\item Self-propagate across networks without a host program
\item Only affect monitors
\end{enumerate}
\answer{C}{Self-propagate across networks without a host program. Worms are self-contained and can spread automatically.}
\source{JKSSB official syllabus plus current-topic extension; original synthesis}
\section*{Question 9}
Ransomware typically:
\begin{enumerate}[label=\Alph*.]
\item Resolves DNS
\item Improves backups
\item Creates a spreadsheet chart
\item Encrypts or blocks access and demands payment
\end{enumerate}
\answer{D}{Encrypts or blocks access and demands payment. Ransomware extorts victims by denying access to data or systems.}
\source{JKSSB official syllabus plus current-topic extension; original synthesis}
\section*{Question 10}
MFA improves authentication by requiring:
\begin{enumerate}[label=\Alph*.]
\item More than one factor
\item Only a longer username
\item A larger monitor
\item No password or factor
\end{enumerate}
\answer{A}{More than one factor. Multi-factor authentication combines independent evidence types.}
\source{JKSSB official syllabus plus current-topic extension; original synthesis}
\section*{Question 11}
The CIA triad stands for:
\begin{enumerate}[label=\Alph*.]
\item Control, Internet, Access
\item Confidentiality, Integrity, Availability
\item Code, Input, Arithmetic
\item Cloud, Interface, Algorithm
\end{enumerate}
\answer{B}{Confidentiality, Integrity, Availability. CIA is a core information-security model.}
\source{JKSSB official syllabus plus current-topic extension; original synthesis}
\section*{Question 12}
A vulnerability is:
\begin{enumerate}[label=\Alph*.]
\item A backup copy
\item The exploit itself always
\item A weakness that may be exploited
\item A network cable
\end{enumerate}
\answer{C}{A weakness that may be exploited. A vulnerability is a weakness; an exploit uses it.}
\source{JKSSB official syllabus plus current-topic extension; original synthesis}
\section*{Question 13}
Open-source software generally makes available:
\begin{enumerate}[label=\Alph*.]
\item A physical motherboard
\item Only screenshots
\item A secret binary with no license
\item Source code under a license permitting defined use/modification
\end{enumerate}
\answer{D}{Source code under a license permitting defined use/modification. Open-source licensing governs access and rights to source code.}
\source{JKSSB official syllabus/question-paper pattern; original synthesis}
\section*{Question 14}
BIOS or UEFI is used during:
\begin{enumerate}[label=\Alph*.]
\item Firmware-level startup and hardware initialization
\item Spreadsheet calculation
\item Email delivery
\item PDF annotation only
\end{enumerate}
\answer{A}{Firmware-level startup and hardware initialization. Firmware initializes hardware and helps begin OS loading.}
\source{JKSSB official syllabus plus current-topic extension; original synthesis}
\section*{Question 15}
E-governance refers to:
\begin{enumerate}[label=\Alph*.]
\item A CPU arithmetic operation
\item Use of ICT to deliver or improve government services
\item A printer ink type
\item A file extension
\end{enumerate}
\answer{B}{Use of ICT to deliver or improve government services. E-governance uses digital systems for public administration/services.}
\source{JKSSB official syllabus/question-paper pattern; original synthesis}
\section*{Question 16}
G2C means:
\begin{enumerate}[label=\Alph*.]
\item General Code Compiler
\item Gateway to Cloud
\item Government to Citizen
\item Government to CPU
\end{enumerate}
\answer{C}{Government to Citizen. G2C describes government services for citizens.}
\source{JKSSB official syllabus/question-paper pattern; original synthesis}
\section*{Question 17}
If a network has 4 subnets with 30 usable host addresses each, total usable host addresses are:
\begin{enumerate}[label=\Alph*.]
\item 90
\item 34
\item 124
\item 120
\end{enumerate}
\answer{D}{120. 4 x 30 = 120 usable addresses. Numerical/application item.}
\source{Sanfoundry, Computer Network chapterwise MCQs; original synthesis}
\section*{Question 18}
BGP is primarily used for:
\begin{enumerate}[label=\Alph*.]
\item Routing between autonomous systems
\item Formatting documents
\item Managing local keyboard input
\item Compressing MP3 files
\end{enumerate}
\answer{A}{Routing between autonomous systems. BGP exchanges reachability information between autonomous systems.}
\source{JKSSB official syllabus plus current-topic extension; original synthesis}
\section*{Question 19}
IPv6 addresses are:
\begin{enumerate}[label=\Alph*.]
\item 32 bits long
\item 128 bits long
\item 16 bits long
\item 8 bits long
\end{enumerate}
\answer{B}{128 bits long. IPv6 uses 128-bit addresses. Numerical/application item.}
\source{JKSSB official syllabus plus current-topic extension; original synthesis}
\section*{Question 20}
A backup is different from RAID because RAID primarily:
\begin{enumerate}[label=\Alph*.]
\item Only stores paper
\item Always encrypts files
\item Improves availability/performance, not an independent backup by itself
\item Replaces authentication
\end{enumerate}
\answer{C}{Improves availability/performance, not an independent backup by itself. RAID redundancy does not protect against all deletion, corruption or ransomware.}
\source{JKSSB official syllabus plus current-topic extension; original synthesis}

\chapter{Programming, Algorithms, Data Structures and Integrated Topics}
\textit{20 questions; numerical/application items are marked in the answer explanation.}
\section*{Question 1}
An algorithm is:
\begin{enumerate}[label=\Alph*.]
\item A finite step-by-step procedure for solving a problem
\item A monitor type
\item A file extension only
\item A network cable
\end{enumerate}
\answer{A}{A finite step-by-step procedure for solving a problem. Algorithms specify ordered steps toward a result.}
\source{Sanfoundry, Computer Fundamentals chapterwise MCQs; original synthesis}
\section*{Question 2}
A flowchart represents an algorithm using:
\begin{enumerate}[label=\Alph*.]
\item Only audio samples
\item Standard graphical symbols
\item Database passwords
\item IP packets
\end{enumerate}
\answer{B}{Standard graphical symbols. Flowcharts visually represent control flow.}
\source{Sanfoundry, Computer Fundamentals chapterwise MCQs; original synthesis}
\section*{Question 3}
A compiler generally translates:
\begin{enumerate}[label=\Alph*.]
\item A PDF into a printer
\item Only one keystroke into a pixel
\item A complete source program into target/object code
\item A MAC address into RAM
\end{enumerate}
\answer{C}{A complete source program into target/object code. A compiler translates source code before execution, subject to language/toolchain.}
\source{Sanfoundry, Computer Fundamentals chapterwise MCQs; original synthesis}
\section*{Question 4}
An assembler translates:
\begin{enumerate}[label=\Alph*.]
\item A spreadsheet into RAM
\item HTML into a photograph
\item Email into DNS
\item Assembly language into machine/object code
\end{enumerate}
\answer{D}{Assembly language into machine/object code. Assembly mnemonics are converted to machine instructions.}
\source{Sanfoundry, Computer Fundamentals chapterwise MCQs; original synthesis}
\section*{Question 5}
Machine language consists of:
\begin{enumerate}[label=\Alph*.]
\item Instructions represented in binary/machine codes
\item Only English sentences
\item Printed symbols only
\item Audio waveforms
\end{enumerate}
\answer{A}{Instructions represented in binary/machine codes. Machine language is directly understood by the processor architecture.}
\source{Sanfoundry, Computer Fundamentals chapterwise MCQs; original synthesis}
\section*{Question 6}
Debugging means:
\begin{enumerate}[label=\Alph*.]
\item Backing up a hard disk only
\item Finding and correcting program faults
\item Creating a slide transition
\item Sending an email
\end{enumerate}
\answer{B}{Finding and correcting program faults. Debugging identifies and fixes defects.}
\source{Sanfoundry, Computer Fundamentals chapterwise MCQs; original synthesis}
\section*{Question 7}
OOP commonly organizes programs around:
\begin{enumerate}[label=\Alph*.]
\item DNS records
\item Only printer ports
\item Objects containing data and behaviour
\item File extensions
\end{enumerate}
\answer{C}{Objects containing data and behaviour. Objects combine state/data and operations/methods conceptually.}
\source{Sanfoundry, Computer Fundamentals chapterwise MCQs; original synthesis}
\section*{Question 8}
A stack follows which principle?
\begin{enumerate}[label=\Alph*.]
\item Broadcast-only
\item FIFO
\item Random-only
\item LIFO
\end{enumerate}
\answer{D}{LIFO. The last item pushed is the first item popped.}
\source{Sanfoundry, Computer Fundamentals chapterwise MCQs; original synthesis}
\section*{Question 9}
A queue normally follows:
\begin{enumerate}[label=\Alph*.]
\item FIFO
\item LIFO
\item Last-address-first only
\item No ordering
\end{enumerate}
\answer{A}{FIFO. The first item entering a queue is normally first to leave.}
\source{Sanfoundry, Computer Fundamentals chapterwise MCQs; original synthesis}
\section*{Question 10}
Which operation adds an item to a stack?
\begin{enumerate}[label=\Alph*.]
\item Pop
\item Push
\item Peek-only
\item Route
\end{enumerate}
\answer{B}{Push. Push inserts an item; pop removes the top item.}
\source{Sanfoundry, Computer Fundamentals chapterwise MCQs; original synthesis}
\section*{Question 11}
Which operation removes the top stack item?
\begin{enumerate}[label=\Alph*.]
\item Compile
\item Push
\item Pop
\item Mount
\end{enumerate}
\answer{C}{Pop. Pop removes the most recently pushed item.}
\source{Sanfoundry, Computer Fundamentals chapterwise MCQs; original synthesis}
\section*{Question 12}
A syntax error violates:
\begin{enumerate}[label=\Alph*.]
\item A file backup rule
\item A monitor resolution rule
\item A network route
\item The grammar rules of a programming language
\end{enumerate}
\answer{D}{The grammar rules of a programming language. Syntax describes how valid program statements are formed.}
\source{Sanfoundry, Computer Fundamentals chapterwise MCQs; original synthesis}
\section*{Question 13}
An interpreter usually executes:
\begin{enumerate}[label=\Alph*.]
\item Translated statements during interpretation
\item Only machine hardware
\item Only printed documents
\item Only database keys
\end{enumerate}
\answer{A}{Translated statements during interpretation. Interpreters translate and execute progressively.}
\source{Sanfoundry, Computer Fundamentals chapterwise MCQs; original synthesis}
\section*{Question 14}
If a stack receives PUSH 4, PUSH 7, then POP, the returned value is:
\begin{enumerate}[label=\Alph*.]
\item 4
\item 7
\item 11
\item 0
\end{enumerate}
\answer{B}{7. 7 is on top and is removed first. Numerical/application item.}
\source{Sanfoundry, Computer Fundamentals chapterwise MCQs; original synthesis}
\section*{Question 15}
A linear search examines items:
\begin{enumerate}[label=\Alph*.]
\item Only through a router
\item Only by binary splitting
\item Sequentially until a match or end is found
\item Only in reverse always
\end{enumerate}
\answer{C}{Sequentially until a match or end is found. Linear search checks elements one after another.}
\source{Sanfoundry, Computer Fundamentals chapterwise MCQs; original synthesis}
\section*{Question 16}
A data structure is used to:
\begin{enumerate}[label=\Alph*.]
\item Create a monitor cable
\item Increase screen brightness
\item Send SMTP messages
\item Organize data for efficient access and modification
\end{enumerate}
\answer{D}{Organize data for efficient access and modification. Data structures provide organized representations for operations.}
\source{Sanfoundry, Computer Fundamentals chapterwise MCQs; original synthesis}
\section*{Question 17}
A program counter stores:
\begin{enumerate}[label=\Alph*.]
\item The address of the next instruction
\item The current monitor pixel
\item The email password
\item The database report
\end{enumerate}
\answer{A}{The address of the next instruction. This register tracks instruction sequencing.}
\source{JKSSB official syllabus plus current-topic extension; original synthesis}
\section*{Question 18}
Two’s complement is commonly used to represent:
\begin{enumerate}[label=\Alph*.]
\item Only RGB colours
\item Signed binary integers
\item Email headers
\item Printer paper sizes
\end{enumerate}
\answer{B}{Signed binary integers. Two’s complement encodes positive and negative integers.}
\source{JKSSB official syllabus plus current-topic extension; original synthesis}
\section*{Question 19}
In a 4-bit unsigned binary number, the maximum value is:
\begin{enumerate}[label=\Alph*.]
\item 8
\item 4
\item 15
\item 16
\end{enumerate}
\answer{C}{15. The range is 0 through 2\^{}4-1 = 15. Numerical/application item.}
\source{JKSSB official syllabus plus current-topic extension; original synthesis}
\section*{Question 20}
The best first response to a negative-stem MCQ is to:
\begin{enumerate}[label=\Alph*.]
\item Choose option C always
\item Ignore the stem
\item Select the longest option automatically
\item Identify words such as NOT/EXCEPT and eliminate carefully
\end{enumerate}
\answer{D}{Identify words such as NOT/EXCEPT and eliminate carefully. Negative wording changes the required selection and must be marked mentally.}
\source{JKSSB official syllabus/question-paper pattern; original synthesis}

\backmatter
\chapter*{Final revision checklist}
\begin{itemize}
\item Revise definitions, contrasts, abbreviations and shortcut near-neighbours.
\item Practise the actual Office version specified in the latest notification.
\item Verify current cyber-law, e-governance and cloud questions against official sources.
\item Solve the official JKSSB papers from the source register under timed conditions.
\item Maintain an error log: topic, wrong option, correct rule and reason for the mistake.
\end{itemize}
\end{document}